\section{Representación del entorno, memoria y mapas}
\label{sec:entorno-mapas}

\subsection{Representación externa del laberinto (\texttt{DEC-01})}
\label{subsec:repr-externa}

Tanto Python (\archivo{pc/mapa.py}) como el carro (\archivo{arduino/agente/src/planner/mapa.cpp}) leen exactamente el mismo archivo de texto~\cite{arquitectura2026doc}. Un mapa de $n$ filas por $m$ columnas tiene dos líneas de cabecera y una grilla de $(2n+1)\times(2m+1)$ caracteres entre \cmd{0} y \cmd{1}. Las cabeceras \cmd{inicio f c N|E|S|O} y \cmd{meta f c} definen salida y llegada en coordenadas de celda, con la fila $0$ arriba (Norte) y la columna $0$ a la izquierda (Oeste). Cualquier texto tras \cmd{\#} se ignora, igual que los espacios dentro de cada fila. En la grilla, la celda $(f, c)$ cae en la posición impar-impar $(2f+1, 2c+1)$; las posiciones con una coordenada par y otra impar guardan las paredes entre celdas (\cmd{1} = pared, \cmd{0} = paso libre), y las posiciones par-par $(2f, 2c)$ son los postes de las esquinas.

\begin{decision}{DEC-01: Formato del mapa en grilla de 0/1}
Al definir el formato discutimos dos opciones: escribir directamente una matriz de enteros del $0$ al $15$ con las máscaras de bits por celda, o usar la grilla explícita $(2n+1)\times(2m+1)$ de \cmd{0} y \cmd{1}. Elegimos la grilla por una razón práctica: se puede leer y editar a mano viendo la forma del laberinto, y hace imposible equivocarse poniendo pared al Este de una celda pero olvidando la pared al Oeste de su vecina. Además, evita pasos de conversión porque el ESP32-S3 y el PC leen el mismo archivo~\cite{arquitectura2026doc}.
\end{decision}

Para que ningún mapa mal formado pase en un lado y falle en el otro, \cmd{desde\_grilla} en \archivo{pc/mapa.py} y \cmd{LectorMapa} en \archivo{mapa.cpp} aplican exactamente las mismas cinco comprobaciones:
\begin{enumerate}[leftmargin=1.4em,nosep]
  \item Todas las filas deben tener igual ancho; alto y ancho deben ser impares y $\ge 3$.
  \item Todo el borde exterior ($y \in \{0, \text{alto}-1\}$ o $x \in \{0, \text{ancho}-1\}$) debe estar cerrado con \cmd{1}.
  \item Ninguna celda $(2f+1, 2c+1)$ puede contener un \cmd{1}.
  \item Deben existir \cmd{inicio} y \cmd{meta} dentro de $[0, n-1]\times[0, m-1]$, con orientación en \{\cmd{N}, \cmd{E}, \cmd{S}, \cmd{O}\}.
  \item En el ESP32-S3, \cmd{MAX\_FILAS = MAX\_COLS = 16} (\archivo{mapa.h}) limita el tamaño a $16\times 16$ celdas (grilla máxima de $33\times 33$).
\end{enumerate}

\subsection{Representación interna, memoria estática y persistencia}
\label{subsec:repr-interna-memoria}

Una vez validada la grilla textual, los dos lectores la convierten a una matriz compacta donde cada celda $(f, c)$ guarda sus paredes en cuatro bits:
\[
  \text{Norte } (\cmd{N}) = 1 = 2^0, \quad
  \text{Este } (\cmd{E}) = 2 = 2^1, \quad
  \text{Sur } (\cmd{S}) = 4 = 2^2, \quad
  \text{Oeste } (\cmd{O}) = 8 = 2^3.
\]
En C++, \cmd{planner::Mapa} guarda esa matriz en \cmd{uint8\_t paredes[16][16]} (solo \SI{256}{bytes} de SRAM), y saber si hay pared frente al carro es una operación de un ciclo: \cmd{paredes[f][c] \& (1 \textless\textless{} ori)} (\cref{fig:grilla-bitmask}).

\begin{figure}[H]
\centering
\begin{tikzpicture}[
  font=\sffamily\footnotesize,
  caja/.style={
    rectangle, rounded corners=3pt, draw=azuldoc, fill=white,
    thick, inner sep=5pt, align=center
  },
  paso/.style={-{Stealth[length=2.3mm,width=1.7mm]}, very thick, color=azuldoc},
  etq/.style={font=\sffamily\scriptsize\bfseries, color=azuldoc}
]

% ============================================================================
% 1. Archivo de texto (2n+1)x(2m+1)
% ============================================================================
\node[caja, fill=codigofondo, anchor=north west] (txt) at (0,0) {
  \textbf{1. Archivo \texttt{.txt} $(2n{+}1)\times(2m{+}1)$}\\[3pt]
  {\ttfamily\scriptsize
  \begin{tabular}{l}
  inicio 0 0 E \\
  meta 1 1 \\
  \textcolor{azuldoc}{\textbf{1}}\textcolor{rojoborde}{1}\textcolor{azuldoc}{\textbf{1}}\textcolor{rojoborde}{1}\textcolor{azuldoc}{\textbf{1}} \\
  \textcolor{rojoborde}{1}\textcolor{unalverde}{\textbf{0}}0\textcolor{unalverde}{\textbf{0}}\textcolor{rojoborde}{1} \\
  \textcolor{azuldoc}{\textbf{1}}\textcolor{rojoborde}{1}\textcolor{azuldoc}{\textbf{1}}0\textcolor{azuldoc}{\textbf{1}} \\
  \textcolor{rojoborde}{1}\textcolor{unalverde}{\textbf{0}}0\textcolor{unalverde}{\textbf{0}}\textcolor{rojoborde}{1} \\
  \textcolor{azuldoc}{\textbf{1}}\textcolor{rojoborde}{1}\textcolor{azuldoc}{\textbf{1}}\textcolor{rojoborde}{1}\textcolor{azuldoc}{\textbf{1}}
  \end{tabular}}\\[3pt]
  {\scriptsize \textcolor{unalverde}{\textbf{0}}: celda $(2f{+}1,2c{+}1)$ \;\; \textcolor{rojoborde}{1}: pared \;\; \textcolor{azuldoc}{\textbf{1}}: poste}
};

% ============================================================================
% 2. Grilla lógica de celdas (f,c) sin tikzpicture anidado
% ============================================================================
\begin{scope}[shift={(5.4,-0.35)}]
  \node[font=\sffamily\footnotesize\bfseries, color=azuldoc] at (1.1,0.3) {2. Celdas lógicas $(f,c)$};
  \fill[azulfondo] (0,-2.0) rectangle (2.2,0.0);
  \draw[grisborde, thin, dashed] (0,-1.0) -- (2.2,-1.0) (1.1,0.0) -- (1.1,-2.0);
  \draw[line width=2pt, color=azuldoc] (0,0) -- (2.2,0) -- (2.2,-2.0) -- (0,-2.0) -- cycle;
  \draw[line width=2pt, color=azuldoc] (0,-1.0) -- (1.1,-1.0);
  \node[font=\ttfamily\scriptsize] at (0.55,-0.5) {(0,0) S};
  \node[font=\ttfamily\scriptsize] at (1.65,-0.5) {(0,1)};
  \node[font=\ttfamily\scriptsize] at (0.55,-1.5) {(1,0)};
  \node[font=\ttfamily\scriptsize, color=unalverde] at (1.65,-1.5) {\textbf{(1,1) M}};
  \node[font=\sffamily\scriptsize, color=gristexto] at (1.1,-2.35) {$f\in[0,n{-}1]$, $c\in[0,m{-}1]$};
  \coordinate (celdas_w) at (-0.15,-1.0);
  \coordinate (celdas_e) at (2.35,-1.0);
\end{scope}

% ============================================================================
% 3. Máscara de 4 bits (uint8_t)
% ============================================================================
\node[caja, fill=verdefondo, draw=unalverde, anchor=north west] (mask) at (8.8,0) {
  \textbf{3. Bitmask \texttt{paredes[f][c]}}\\[2pt]
  {\scriptsize \texttt{N=1, E=2, S=4, O=8}}\\[4pt]
  {\ttfamily\scriptsize
  \begin{tabular}{|c|c|}
  \hline
  \textbf{13} (\texttt{1101}) & \textbf{3} (\texttt{0011}) \\
  {\tiny N+S+O} & {\tiny N+E} \\
  \hline
  \textbf{13} (\texttt{1101}) & \textbf{6} (\texttt{0110}) \\
  {\tiny N+S+O} & {\tiny E+S} \\
  \hline
  \end{tabular}}\\[4pt]
  {\scriptsize Acceso $O(1)$: \texttt{paredes[f][c] \& (1 \textless\textless{} ori)}}
};

% ============================================================================
% 4. Plano físico en cm (celda 24 cm, línea 2 cm) sin tikzpicture anidado
% ============================================================================
\begin{scope}[shift={(4.6,-6.2)}, scale=0.065]
  \fill[ambarfondo, rounded corners=3pt] (-14,-14) rectangle (66,64);
  \draw[ambarborde, thick, rounded corners=3pt] (-14,-14) rectangle (66,64);
  \node[font=\sffamily\footnotesize\bfseries, color=azuldoc] at (26,58)
    {4. Plano físico en pista (\texttt{pc/imprimir.py}, en cm)};
  % Margen 5 cm
  \fill[white] (-5,-5) rectangle (55,55);
  \draw[grisborde, dashed] (-5,-5) rectangle (55,55);
  % Paredes negras de grosor 2 cm (celda = 24 cm -> 50x50 cm)
  \fill[black] (0,48) rectangle (50,50);   % H1 arriba
  \fill[black] (0,24) rectangle (26,26);   % H2 media izq
  \fill[black] (0,0) rectangle (50,2);     % H3 abajo
  \fill[black] (0,2) rectangle (2,24);     % V izq inf
  \fill[black] (0,26) rectangle (2,48);    % V izq sup
  \fill[black] (48,2) rectangle (50,48);   % V der
  % Cotas
  \draw[{Stealth[length=1.5mm]}-{Stealth[length=1.5mm]}, color=azuldoc, thick]
    (2,13) -- (26,13) node[midway, above, font=\sffamily\tiny] {22 cm libre};
  \draw[{Stealth[length=1.5mm]}-{Stealth[length=1.5mm]}, color=rojoborde, thick]
    (1,37) -- (25,37) node[midway, above, font=\sffamily\tiny] {24 cm paso};
  \node[font=\sffamily\tiny, color=gristexto] at (26,-10)
    {Ancho total $= 24c + 2 + 2\cdot 5 = \SI{60}{cm}$ (cinta mate / cartulina)};
  \coordinate (fisico_e) at (66,25);
\end{scope}

% Flechas de flujo
\draw[paso] (txt.east) -- (celdas_w);
\draw[paso] (celdas_e) -- (mask.west);
\draw[paso] (mask.south) |- node[pos=0.25, right, etq] {escala 1:1} (fisico_e);

\end{tikzpicture}

\caption{Transformación de la representación del laberinto: archivo de texto $(2n+1)\times(2m+1)$ $\rightarrow$ grilla lógica de celdas $(f,c)$ $\rightarrow$ matriz de máscaras de 4 bits $\rightarrow$ plano físico en centímetros.}
\label{fig:grilla-bitmask}
\end{figure}

\paragraph{Memoria estática y persistencia en el ESP32-S3.}
En un microcontrolador conviene evitar el \textit{heap} y tampoco llenar el \textit{stack} del \cmd{loop()}. Por eso todos los arreglos de búsqueda son estáticos y acotados a $\text{MAX\_ESTADOS} = 16 \times 16 \times 4 = 1024$. En \archivo{bfs.cpp} (en \texttt{main}), los cuatro arreglos \cmd{padre}, \cmd{accion}, \cmd{visto} y \cmd{cola} ocupan apenas $\approx\SI{4}{KB}$ de SRAM. En \archivo{astar.cpp} (rama \texttt{firmware/astar}), la frontera usa un montículo binario estático de $3\cdot\text{MAX\_ESTADOS}+1 = 3073$ nodos (\cmd{double f}, \cmd{uint16\_t orden}, \cmd{int16\_t estado}) junto con \cmd{g[1024]} en \texttt{double} para empatar exactamente igual que Python, lo que suma $\approx\SI{60}{KB}$ de SRAM fija.

Para no perder el mapa ni la calibración si se apaga el carro o entra un \textit{brownout}, el firmware guarda todo en Flash no volátil (NVS) mediante \cmd{Preferences} bajo el \textit{namespace} \cmd{"agente"}: el bloque binario del mapa (\cmd{prefs.putBytes("mapa", \&mapa, sizeof(mapa))}), \cmd{costo\_giro}, los parámetros de \cmd{SET}/\cmd{SAVE} y las credenciales del \textit{hotspot} (\cmd{red\_ssid}, \cmd{red\_clave})~\cite{arquitectura2026doc}.

\paragraph{Alcance de la actualización del mapa.}
En esta versión el mapa se conoce de antemano y no se modifica durante la corrida: el carro no hace SLAM ni descubre paredes nuevas. Mientras ejecuta \cmd{GO}, solo actualiza su pose discreta \cmd{ej.pose} tras cada tramo y el rumbo continuo en \cmd{imu::rumbo()}. De todas formas, como \cmd{paredes[f][c]} solo usa 4 de los 8 bits de cada \cmd{uint8\_t}, los 4 bits altos quedan libres si más adelante queremos marcar celdas visitadas o paredes descubiertas en línea.

\subsection{Mapas de prueba 2D, generador aleatorio e impresión física}
\label{subsec:mapas-generacion-impresion}

La \cref{tab:mapas-repo} lista los siete mapas creados en el repositorio, separando los de simulación de las pistas físicas pensadas para el carro.

\begin{table}[H]
\centering
\caption{Inventario de mapas de prueba 2D en el repositorio y ramas activas.}
\label{tab:mapas-repo}
\small
\begin{tabularx}{\textwidth}{@{}p{3.4cm}p{1.5cm}p{2.2cm}p{2.3cm}X@{}}
\toprule
\textbf{Archivo} & \textbf{Tamaño} & \textbf{Inicio $\rightarrow$ Meta} & \textbf{Estado} & \textbf{Propósito y características} \\
\midrule
\archivo{prueba\_4x4.txt} & $4\times 4$ & \estado{0}{0}{\cmd{E}} $\rightarrow$ $(3,3)$ & \badgeok & Primera pista impresa (celda \SI{22}{cm}, margen \SI{0.5}{cm}, $91\times 91\text{ cm}$) para pruebas del 26/10/2026. \\
\archivo{ejemplo\_5x5.txt} & $5\times 5$ & \estado{0}{0}{\cmd{E}} $\rightarrow$ $(4,4)$ & \badgeok & Mapa de referencia para \cmd{planear.py}, \cmd{enviar.py} y pruebas unitarias. \\
\archivo{simulado\_6x6.txt} & $6\times 6$ & \estado{0}{0}{\cmd{E}} $\rightarrow$ $(5,5)$ & \badgepr{} (\texttt{mapas/SW-05}) & Laberinto aleatorio con ciclos ($p=0.2$, semilla 106). \\
\archivo{simulado\_10x10.txt} & $10\times 10$ & \estado{0}{0}{\cmd{E}} $\rightarrow$ $(9,9)$ & \badgepr{} (\texttt{mapas/SW-05}) & Evaluación intermedia de escala ($p=0.2$, semilla 38). \\
\archivo{simulado\_16x16.txt} & $16\times 16$ & \estado{0}{0}{\cmd{E}} $\rightarrow$ $(15,15)$ & \badgepr{} (\texttt{mapas/SW-05}) & Prueba en el límite máximo de memoria del ESP32-S3 ($p=0.2$, semilla 23). \\
\archivo{fisicos/mapa1.txt} & $4\times 3$ & \estado{0}{0}{\cmd{E}} $\rightarrow$ $(3,2)$ & \badgepr{} (\texttt{mapas/mapa2}) & Pista física (\SI{24}{cm}/celda, $84\times 108\text{ cm}$ con margen) con paredes al final de los tramos rectos para alinear con el frontal. \\
\archivo{fisicos/mapa2.txt} & $4\times 3$ & \estado{0}{0}{\cmd{E}} $\rightarrow$ $(3,2)$ & \badgepr{} (\texttt{mapas/mapa2}) & Pista física (\SI{24}{cm}/celda) con intersección abierta: el robot debe girar aunque tenga paso libre al frente. \\
\bottomrule
\end{tabularx}
\end{table}

Para comparar algoritmos en laberintos con varios caminos alternativos, escribimos \archivo{pc/generar.py} en la rama \texttt{mapas/SW-05}. El script crea mapas válidos entre $1\times 1$ y $16\times 16$ en dos pasos: primero talla un árbol generador sin ciclos con un DFS iterativo (lo que asegura que siempre exista al menos una ruta entre inicio y meta), y después recorre las paredes internas sobrantes abriéndolas con probabilidad \cmd{ciclos} $\in [0, 1]$ bajo una semilla fija (\cmd{random.Random(semilla)}) para que cualquier experimento se pueda repetir.

Del lado físico, \archivo{pc/imprimir.py} (ya en \texttt{main}) convierte cualquier archivo de mapa en dos PDF dentro de \archivo{build/impresion/}. Primero \cmd{piezas()} junta las paredes horizontales colineales en tiras enteras (\cmd{H1, H2, ...}) y recorta las verticales (\cmd{V1, V2, ...}) donde tocan una horizontal para no encimar dos capas de cinta o cartulina. Luego \cmd{tiras\_necesarias()} acomoda esos cortes con una heurística \textit{First-Fit Decreasing} para calcular cuántos pliegos de cartulina o metros de cinta se necesitan. Las dos salidas son \archivo{<mapa>\_plotter.pdf} (escala 1:1 con una regla de \SI{10}{cm} en el borde para comprobar que la impresora no encogió el plano) y \archivo{<mapa>\_armado.pdf} (cotas y lista de corte para armar el laberinto a mano).

\begin{hallazgo}[Invisibilidad infrarroja de la tinta de plotter (26/10/2026)]
Al probar el primer plano impreso de \archivo{prueba\_4x4.txt}, los TCRT5000 (\SI{950}{nm}) leyeron la tinta negra del plotter casi igual que el papel blanco ($\approx 24\text{--}85$ frente a $400\text{--}690$ sobre cinta aislante negra mate)~\cite{arquitectura2026doc}. Por eso usamos el plano impreso solo como plantilla de medidas y forramos todas las líneas con cinta aislante negra mate de $\approx\SI{1.9}{cm}$ centrada sobre el trazo de \SI{2}{cm}.
\end{hallazgo}
